\chapter{Tools, Skills e MCP}
\label{chap:tools_skills_mcp}

% ---------------------------------------------------------------
\section{Pr\'e-requisitos}
% ---------------------------------------------------------------

Para aproveitar este cap\'itulo, o leitor deve ter conclu\'ido:

\begin{itemize}[leftmargin=*]
  \item Cap\'itulo~\ref{chap:agentes_langgraph_memoria} --- ciclo ReAct
        e grafos de estado com LangGraph;
  \item Familiaridade com fun\c{c}\~oes ass\'incronas em Python
        (\texttt{async}/\texttt{await}).
\end{itemize}

% ---------------------------------------------------------------
\section{Escopo do cap\'itulo}
% ---------------------------------------------------------------

\begin{quote}\itshape
\textbf{Onde estamos na hist\'oria:} loop interno, segunda virada.
O agente j\'a sabe raciocinar em ciclos (cap.~3). Agora ele precisa
de m\~aos: ferramentas que consultam bancos de dados, APIs e fontes
externas. O protocolo MCP \'e a linguagem padronizada com que o
agente declara e invoca essas ferramentas --- a mesma linguagem que
conectar\'a, nos cap\'itulos 7--10, ao pipeline de produ\c{c}\~ao.
\end{quote}

Este cap\'itulo apresenta tr\^es camadas de extensibilidade do agente:

\begin{enumerate}[leftmargin=*]
  \item \textbf{Tools nativas LangChain:} fun\c{c}\~oes Python decoradas
        com \texttt{@tool} que o agente pode invocar diretamente;
  \item \textbf{Skills compostas:} agrupamento de tools com contexto
        de neg\'ocio embutido (ex.\ \textit{skill} de investimentos);
  \item \textbf{MCP (\textit{Model Context Protocol}):} protocolo
        padronizado pela Anthropic que expõe ferramentas via servidor
        \texttt{stdio} ou HTTP, permitindo reutiliza\c{c}\~ao entre
        agentes e equipes. No AWS Prescriptive Guidance, esta camada
        corresponde ao padr\~ao de \textit{tool-based agents for servers}~\cite{awsagentic2025}.
\end{enumerate}

% ---------------------------------------------------------------
\section{Objetivos de aprendizagem}
% ---------------------------------------------------------------

Ao concluir este cap\'itulo, o leitor ser\'a capaz de:

\begin{itemize}[leftmargin=*]
  \item Criar e registrar tools nativas no grafo LangGraph;
  \item Diferenciar tool, skill e servidor MCP;
  \item Implementar um servidor MCP com transporte \texttt{stdio}
        usando a biblioteca \texttt{mcp};
  \item Conectar um agente LangGraph a um servidor MCP externo;
  \item Validar a integra\c{c}\~ao com casos de teste offline.
\end{itemize}

% ---------------------------------------------------------------
\section{Conceitos te\'oricos}
% ---------------------------------------------------------------

\subsection{Tools nativas: fun\c{c}\~oes como a\c{c}\~oes do agente}

No LangGraph\index{LangGraph}, uma \textit{tool}\index{tool (agente)|textbf}
\'e qualquer fun\c{c}\~ao Python decorada com \texttt{@tool}
do LangChain. O agente recebe a lista de tools no momento da
constru\c{c}\~ao do grafo e decide, a cada n\'o de racioc\'inio
ReAct\index{ReAct}, qual tool invocar e com quais argumentos.

\begin{lstlisting}[language=Python,
  caption={Tool nativa: consulta de saldo de cliente.},
  label={lst:tool_nativa}]
from langchain_core.tools import tool

@tool
def consultar_saldo(cliente_id: str) -> str:
    """Retorna o saldo atual do cliente. Use quando o usuario perguntar
    sobre saldo, limite ou posicao financeira."""
    # Em producao: chamada ao DynamoDB ou API core bancario
    saldos = {"C001": 15000.00, "C002": 8500.50, "C003": 22000.00}
    saldo = saldos.get(cliente_id)
    if saldo is None:
        return f"Cliente {cliente_id} nao encontrado."
    return f"Saldo do cliente {cliente_id}: R$ {saldo:,.2f}"
\end{lstlisting}

O \textit{docstring} da fun\c{c}\~ao \'e o que o LLM l\^e para
decidir quando usar a tool. Docstrings claros e com exemplos de
contexto de uso melhoram significativamente a precis\~ao do agente.

\subsection{Skills: tools com contexto de neg\'ocio}

Uma \textit{skill}\index{skill (agente)|textbf} \'e um conjunto de
tools relacionadas que formam uma capacidade de neg\'ocio coesa.
No projeto banc\'ario, a \textit{skill} de investimentos agrupa as
tools de consulta de produtos, simula\c{c}\~ao de rendimento e
recomenda\c{c}\~ao por perfil. A vantagem \'e que o agente pode
ser instanciado com skills diferentes conforme o segmento do cliente,
sem alterar a l\'ogica do grafo.

\subsection{MCP: protocolo padronizado para ferramentas}
\label{sec:mcp_protocolo}

O \textit{Model Context Protocol}
(MCP)\index{MCP (\textit{Model Context Protocol})|textbf} \'e um
protocolo aberto que padroniza como servidores de ferramentas
comunicam suas capacidades a agentes de IA~\cite{mcp2024}.
A arquitetura tem tr\^es pap\'eis:

\begin{itemize}[leftmargin=*]
  \item \textbf{Cliente MCP:} o agente (LangGraph, Claude Desktop,
        etc.) que solicita ferramentas;
  \item \textbf{Servidor MCP:} processo que expõe ferramentas via
        \texttt{stdio} (padr\~ao para integra\c{c}\~ao local) ou
        HTTP/SSE (para servidores remotos);
  \item \textbf{Host:} orquestrador que gerencia a conex\~ao entre
        cliente e servidor (ex.\ Claude Desktop, VS~Code MCP
        extension).
\end{itemize}

A comunica\c{c}\~ao segue o protocolo JSON-RPC~2.0. O servidor
responde a tr\^es opera\c{c}\~oes principais: \texttt{tools/list}
(anuncia ferramentas), \texttt{tools/call} (executa uma ferramenta)
e \texttt{resources/list} (expõe recursos est\'aticos).

A vantagem do MCP sobre tools nativas \'e o \textbf{desacoplamento}:
o servidor MCP pode ser desenvolvido em qualquer linguagem,
versionado separadamente e reutilizado por m\'ultiplos agentes
e equipes.

% ---------------------------------------------------------------
\section{Implementa\c{c}\~ao orientada por passos}
% ---------------------------------------------------------------

\subsection{Passo 1 --- Criar tools nativas e registrar no agente}

Com base no grafo ReAct do cap\'itulo anterior, adicionamos tools
espec\'ificas ao dominio banc\'ario:

\begin{lstlisting}[language=Python,
  caption={Tools nativas para o agente banc\'ario.},
  label={lst:tools_bancarias}]
from langchain_core.tools import tool
from langgraph.prebuilt import create_react_agent
from langchain_aws import ChatBedrockConverse

@tool
def consultar_saldo(cliente_id: str) -> str:
    """Retorna o saldo atual do cliente."""
    saldos = {"C001": 15000.00, "C002": 8500.50}
    s = saldos.get(cliente_id)
    return f"R$ {s:,.2f}" if s else f"Cliente {cliente_id} nao encontrado."

@tool
def listar_produtos_investimento(perfil: str) -> str:
    """Lista produtos de investimento adequados ao perfil.
    Perfis validos: conservador, moderado, arrojado."""
    catalogo = {
        "conservador": ["Tesouro Selic", "CDB 100% CDI", "LCI isenta IR"],
        "moderado":    ["Fundo Multimercado", "CRA", "Debentures"],
        "arrojado":    ["Acoes growth", "FII", "COE estruturado"],
    }
    produtos = catalogo.get(perfil.lower(), [])
    if not produtos:
        return f"Perfil '{perfil}' nao reconhecido."
    return "Produtos recomendados: " + ", ".join(produtos)

llm = ChatBedrockConverse(
    model="us.anthropic.claude-sonnet-4-5-20250514-v1:0",
    region_name="us-east-1",
)
tools = [consultar_saldo, listar_produtos_investimento]
agente = create_react_agent(llm, tools)
\end{lstlisting}

\subsection{Passo 2 --- Implementar servidor MCP com transporte stdio}

O arquivo \texttt{mcp\_mock\_server.py} implementa um servidor MCP
completo com dados mock, ideal para testes offline. A estrutura
\'e composta por tr\^es partes:

\begin{lstlisting}[language=Python,
  caption={Estrutura do servidor MCP mock (simplificado).},
  label={lst:mcp_server_estrutura}]
import asyncio, json
from mcp.server import Server
from mcp.server.stdio import stdio_server
import mcp.types as types

server = Server("postgres-mock-server")

@server.list_tools()
async def list_tools() -> list[types.Tool]:
    return [
        types.Tool(
            name="executar_query",
            description="Executa query no banco mock.",
            inputSchema={
                "type": "object",
                "properties": {
                    "tabela": {"type": "string"},
                    "filtro": {"type": "object"}
                },
                "required": ["tabela"]
            }
        ),
        types.Tool(name="listar_tabelas",
                   description="Lista tabelas disponiveis.",
                   inputSchema={"type": "object", "properties": {}}),
    ]

@server.call_tool()
async def call_tool(name: str, arguments: dict) -> list[types.TextContent]:
    if name == "listar_tabelas":
        return [types.TextContent(type="text",
                                  text=json.dumps(list(MOCK_DATA.keys())))]
    # ... logica de execar_query

async def main():
    async with stdio_server() as (read, write):
        await server.run(read, write,
                         server.create_initialization_options())

if __name__ == "__main__":
    asyncio.run(main())
\end{lstlisting}

O transporte \texttt{stdio} significa que o agente inicia o servidor
como subprocesso e se comunica via \texttt{stdin}/\texttt{stdout}.
Essa arquitetura \'e \textbf{segura por padr\~ao}: sem portas abertas,
sem autentica\c{c}\~ao de rede necess\'aria para uso local.

\subsection{Passo 3 --- Conectar o agente ao servidor MCP}

O LangGraph suporta conex\~ao a servidores MCP via
\texttt{MultiServerMCPClient}. Em produ\c{c}\~ao, o servidor
\texttt{mcp\_postgres\_server.py} conecta ao PostgreSQL real; em
testes, usa-se o mock:

\begin{lstlisting}[language=Python,
  caption={Conex\~ao do agente ao servidor MCP.},
  label={lst:mcp_client}]
import asyncio
from langchain_mcp_adapters.client import MultiServerMCPClient
from langgraph.prebuilt import create_react_agent
from langchain_aws import ChatBedrockConverse

async def criar_agente_com_mcp():
    async with MultiServerMCPClient({
        "banco_mock": {
            "command": "python",
            "args": ["mcp_mock_server.py"],
            "transport": "stdio",
        }
    }) as client:
        tools = client.get_tools()
        llm = ChatBedrockConverse(
            model="us.anthropic.claude-sonnet-4-5-20250514-v1:0",
            region_name="us-east-1",
        )
        agente = create_react_agent(llm, tools)
        resultado = await agente.ainvoke({
            "messages": [{"role": "user",
                          "content": "Quais tabelas estao disponiveis?"}]
        })
        return resultado["messages"][-1].content

if __name__ == "__main__":
    print(asyncio.run(criar_agente_com_mcp()))
\end{lstlisting}

\subsection{Passo 4 --- Decis\~oes de projeto: tool nativa vs.\ MCP}

A escolha entre tool nativa e servidor MCP deve seguir o crit\'erio
de acoplamento:

\begin{table}[htbp]
\centering
\caption{Crit\'erios de escolha: tool nativa vs.\ MCP}
\label{tab:tool_vs_mcp}
\begin{tabular}{lll}
\toprule
\textbf{Crit\'erio} & \textbf{Tool nativa} & \textbf{Servidor MCP} \\
\midrule
Deploy conjunto com agente & \checkmark & --- \\
Reutilizavel por outros agentes & --- & \checkmark \\
Linguagem diferente de Python & --- & \checkmark \\
Dados sens\'iveis isolados & --- & \checkmark \\
Latencia m\'inima & \checkmark & --- \\
Versionamento independente & --- & \checkmark \\
\bottomrule
\end{tabular}
\end{table}

No projeto banc\'ario, tools de c\'alculo interno (\'indice de
confiabilidade, formata\c{c}\~ao de resposta) s\~ao nativas; acesso
a bases de dados (PostgreSQL, DynamoDB) \'e feito via MCP para
garantir isolamento e auditabilidade.

% ---------------------------------------------------------------
\section{Autoavalia\c{c}\~ao}
% ---------------------------------------------------------------

Marque V (verdadeiro) ou F (falso):

\begin{enumerate}[leftmargin=*]
  \item O \textit{docstring} da fun\c{c}\~ao decorada com
        \texttt{@tool} n\~ao influencia o comportamento do agente. \\
        \textbf{Gabarito: F} --- O \textit{docstring} \'e o texto
        que o LLM l\^e para decidir quando e como usar a tool;
        docstrings imprecisos reduzem a qualidade das decis\~oes.

  \item No transporte \texttt{stdio}, o servidor MCP abre uma porta
        TCP para receber chamadas do agente. \\
        \textbf{Gabarito: F} --- No transporte \texttt{stdio}, a
        comunica\c{c}\~ao ocorre via \texttt{stdin}/\texttt{stdout}
        do subprocesso, sem portas de rede.

  \item Um servidor MCP pode ser implementado em qualquer linguagem
        que suporte JSON-RPC~2.0. \\
        \textbf{Gabarito: V} --- O protocolo \'e independente de
        linguagem; existem SDKs oficiais em Python, TypeScript,
        Kotlin e Rust.

  \item Tools nativas s\~ao sempre prefer\'iveis ao MCP por terem
        menor lat\^encia. \\
        \textbf{Gabarito: F} --- A lat\^encia menor n\~ao justifica
        o acoplamento quando o recurso precisa ser isolado ou
        reutilizado por outros agentes.
\end{enumerate}

% ---------------------------------------------------------------
\section{Resumo}
% ---------------------------------------------------------------

\begin{itemize}[leftmargin=*]
  \item Tools nativas s\~ao fun\c{c}\~oes Python com \texttt{@tool};
        o \textit{docstring} guia a decis\~ao do LLM.
  \item Skills agrupam tools relacionadas para uma capacidade de
        neg\'ocio coesa.
  \item MCP padroniza a comunica\c{c}\~ao entre agentes e servidores
        de ferramentas via JSON-RPC~2.0.
  \item O transporte \texttt{stdio} \'e a op\c{c}\~ao padr\~ao para
        uso local, seguro e sem portas abertas.
  \item A escolha tool nativa vs.\ MCP depende de acoplamento,
        reutiliza\c{c}\~ao e necessidade de isolamento.
  \item No projeto banc\'ario, acesso a dados sensíveis \'e sempre
        mediado por servidores MCP.
\end{itemize}

% ---------------------------------------------------------------
\section{Para saber mais}
% ---------------------------------------------------------------

\begin{itemize}[leftmargin=*]
  \item \citeonline{mcp2024} --- especifica\c{c}\~ao oficial do
        \textit{Model Context Protocol}.
  \item Reposit\'orio \texttt{langchain-mcp-adapters}:
        \footnote{\url{https://github.com/langchain-ai/langchain-mcp-adapters}}
        integra\c{c}\~ao oficial do LangGraph com servidores MCP.
  \item \citeonline{langgraph2024} --- documenta\c{c}\~ao de tools
        e \textit{tool calling} no LangGraph.
\end{itemize}

% ---------------------------------------------------------------
\section{Plano de testes do cap\'itulo}
% ---------------------------------------------------------------

\subsection{Testes funcionais}

\begin{itemize}[leftmargin=*]
  \item \textbf{TF-04-01:} invocar \texttt{consultar\_saldo} com
        \texttt{cliente\_id="C001"} e verificar retorno
        \texttt{"R\$ 15.000,00"}.
  \item \textbf{TF-04-02:} invocar \texttt{listar\_tabelas} via
        servidor MCP mock e verificar que retorna lista com
        \texttt{["clientes","produtos","pedidos"]}.
  \item \textbf{TF-04-03:} invocar \texttt{executar\_query} com
        \texttt{filtro=\{"status":"entregue"\}} e verificar 2 registros.
  \item \textbf{TF-04-04:} criar agente com as duas tools nativas e
        verificar que responde corretamente \`a pergunta
        \textit{``Qual o saldo do cliente C002?''}.
\end{itemize}

\subsection{Teste de isolamento}

\begin{itemize}[leftmargin=*]
  \item \textbf{TI-04-01:} verificar que o servidor MCP n\~ao abre
        nenhuma porta de rede ao rodar com transporte \texttt{stdio}.
  \item \textbf{TI-04-02:} verificar que tool com filtro
        de cliente inexistente retorna mensagem de erro sem
        exce\c{c}\~ao n\~ao tratada.
\end{itemize}

% ---------------------------------------------------------------
\section{Crit\'erios de aceite}
% ---------------------------------------------------------------

\begin{itemize}[leftmargin=*]
  \item \texttt{mcp\_mock\_server.py} responde corretamente \`as tr\^es
        opera\c{c}\~oes: \texttt{listar\_tabelas},
        \texttt{descrever\_tabela} e \texttt{executar\_query}.
  \item Agente com tools nativas invoca a tool correta para cada
        tipo de pergunta (saldo vs.\ produtos).
  \item Todos os testes TF-04-0x e TI-04-0x passam com
        \texttt{python validate\_sprint\_a.py}.
\end{itemize}

% ---------------------------------------------------------------
\section{Espa\c{c}o para expans\~ao iterativa}
% ---------------------------------------------------------------

\begin{itemize}[leftmargin=*]
  \item \textbf{v4.1:} servidor MCP com transporte HTTP/SSE para
        deploy em container ECS.
  \item \textbf{v4.2:} autent\'ica\c{c}\~ao OAuth 2.0 no servidor
        MCP (requerido para acesso a dados PCI-DSS).
  \item \textbf{Lacuna:} benchmark de lat\^encia tool nativa vs.\
        MCP em cen\'ario de alta concorr\^encia.
\end{itemize}
